
\section{$A_\infty$ 代数的同伦等价：几何实现}
\label{sec:4-6}
\liftIndex{I172}

\subsection{$A_\infty$ 同态的构造}
\label{subsec:4-6-1}

本节旨在证明定理 3.5.11 所构造的滤过 $A_\infty$ 代数
不依赖于构造中涉及的各种选择。更确切地说，我们将证明下面的
定理 4.6.1（以及定理 4.1.1）。回忆第 4.1 节的记号。
设 $(M,\omega)$、$(M',\omega')$ 是辛流形，
$L\subset M$、$L'\subset M'$ 是 $n$ 维拉格朗日子流形。
设 $\psi\colon M\to M'$ 是满足 $\psi(L)=L'$ 的辛微分同胚。
令 $J$ 为 $M$ 上的相容殆复结构。我们如第 3.5 节一样，
用 $J$ 构造模空间 $\mathcal M_{k+1}^{\mathrm{main}}(\beta)$；
为标明 $L,J$ 的选择，记为
$\mathcal M_{k+1}^{\mathrm{main}}((L,J),\beta)$。
另取 $M'$ 上的相容殆复结构 $J'$，同样定义
$\mathcal M_{k+1}^{\mathrm{main}}((L',J'),\beta)$。
这里不对 $J,J'$ 假定任何一般性条件；
所需的横截性由抽象扰动而非扰动 $J$ 得到。

在第 3.5 节，我们取了光滑奇异单形的可数集
$\mathcal X_1(L)$，并考察它张成的 $S(L;\mathbb Q)$
的子复形，记为 $C_1(L;\mathbb Q)$。
对 $P_i\in\mathcal X_1(L)$，还考察了纤维积
\[
 \mathcal M_{k+1}^{\mathrm{main}}((L,J),\beta;\vec P)
 =
 \mathcal M_{k+1}^{\mathrm{main}}((L,J),\beta)
 {}_{ev}\!\times(P_1\times\cdots\times P_k).
\]
这是带仓西结构的空间。命题 3.5.2（其证明将在第 7.2 节完成）
断言：可以选择 $\mathcal X_1(L)$ 和
$\mathcal M_{k+1}^{\mathrm{main}}((L,J),\beta;\vec P)$
上的一套多值截面 $\mathfrak s$，使得对任意
$P_i\in\mathcal X_1(L)$，把
$\mathfrak s^{-1}(0)$ 适当三角剖分后，
\[
 \mathcal M_{k+1}^{\mathrm{main}}
 ((L,J),\beta;\vec P)^{\mathfrak s}
 =
 (\mathfrak s^{-1}(0),ev_0)
\]
属于 $C_1(L;\mathbb Q)$。命题 3.5.2 还断言，
这些多值截面在
$\mathcal M_{k+1}^{\mathrm{main}}((L,J),\beta;\vec P)$
的紧化边界上相容。于是定义
\[
 \mathfrak m_{k,\beta}(P_1,\ldots,P_k)
 =
 \mathcal M_{k+1}^{\mathrm{main}}
 ((L,J),\beta;\vec P)^{\mathfrak s}.
\]
为标明 $J,\mathfrak s$ 的选择，写作
$\mathfrak m_{k,\beta}^{J,\mathfrak s}$。
这些映射在 $C_1(L;\Lambda_{0,\mathrm{nov}})$
上定义一个滤过 $A_\infty$ 代数结构。

对 $L',J'$ 作类似选择，得到光滑奇异单形的可数集
$\mathcal X_2(L')$、它张成的子复形
$C_2(L';\mathbb Q)$，以及在 $P_i'\in\mathcal X_2(L')$
时的模空间
$\mathcal M_{k+1}^{\mathrm{main}}
((L',J'),\beta;\vec P')^{\mathfrak s'}$。
这在 $C_2(L';\Lambda_{0,\mathrm{nov}})$
上诱导滤过 $A_\infty$ 代数结构，
记其运算为 $\mathfrak m_{k,\beta}^{J',\mathfrak s'}$。
回忆 $C_1(L;\mathbb Q)$、$C_2(L';\mathbb Q)$
的选取保证其上同调群分别同构于 $L$、$L'$ 的上同调群。

\begin{theorem}{4.6.1}\label{thm:4-6-1}
存在光滑奇异单形的可数集 $\mathcal X_3(L')$，
它张成 $S(L';\mathbb Q)$ 的子复形
$C_3(L';\mathbb Q)$；存在
$C_3(L';\Lambda_{0,\mathrm{nov}})$ 上的滤过
$A_\infty$ 代数结构
$\mathfrak m_{k,\beta}^{J',\mathfrak s''}$；
还存在有隙滤过 $A_\infty$ 同态
\[
 \mathfrak f\colon C_1(L;\Lambda_{0,\mathrm{nov}})
 \longrightarrow C_3(L';\Lambda_{0,\mathrm{nov}}),
\]
并满足以下性质：
\begin{enumerate}[label=\textup{(4.6.2.\arabic*)},leftmargin=5.9em]
 \item
 $\mathfrak s''$ 是
 $\mathcal M_{k+1}^{\mathrm{main}}
 ((L',J'),\beta;\vec P)$
 上的一套多值截面，并在边界处相容；
 这里 $P_i\in C_3(L';\mathbb Q)$。
 \item
 $\mathcal X_3(L')\supset\mathcal X_2(L')$。
 此外，若 $P_i\in\mathcal X_2(L')$，
 则在模空间
 $\mathcal M_{k+1}^{\mathrm{main}}
 ((L',J'),\beta;\vec P)$ 上
 $\mathfrak s''=\mathfrak s'$。
 \item
 $C_3(L';\mathbb Q)$ 关于通常上边界算子的
 上同调同构于 $L'$ 的上同调。
 \item
 $\mathfrak f$ 所诱导的同态
 $(\overline{\mathfrak f}_1)_*=(\psi^{-1})^*$
 在下列辨识下，
 \[
 \begin{gathered}
 H^*(C_1(L;\mathbb Q),\overline{\mathfrak m}_1)
 \cong H^*(L;\mathbb Q)\\
 \xrightarrow{\;(\psi^{-1})^*\;}
 H^*(L';\mathbb Q)
 \cong H^*(C_3(L';\mathbb Q),
 \overline{\mathfrak m}_1)
 \end{gathered}
 \]
 与微分同胚 $\psi^{-1}$ 所诱导的同构
 $(\psi^{-1})^*$ 一致。
\end{enumerate}
\end{theorem}

\begin{corollary}{4.6.3}\label{cor:4-6-3}
滤过 $A_\infty$ 代数
$(C_1(L;\Lambda_{0,\mathrm{nov}}),
\mathfrak m_k^{J,\mathfrak s})$
与
$(C_2(L';\Lambda_{0,\mathrm{nov}}),
\mathfrak m_k^{J',\mathfrak s'})$
同伦等价。
\end{corollary}
推论 4.6.3 是定理 4.1.1 的另一种表述。

\begin{proof}[定理 4.6.1 推出推论 4.6.3]
由 (4.6.2.4)，$\mathfrak f$ 是弱同伦等价；
故由定理 4.2.45，它是同伦等价。另一方面，
由 (4.6.2.2)，自然包含
$C_2(L';\mathbb Q)\subset C_3(L';\mathbb Q)$
诱导一个滤过 $A_\infty$ 同态。
由 (4.6.2.3) 和定理 4.2.45，它也是同伦等价。
再由推论 4.2.44，得到推论 4.6.3。
\end{proof}

本小节其余部分用于证明定理 4.6.1。严格说来，
我们在本小节预先假定命题 4.6.14 成立；
它涉及可数集 $\mathcal X_3(L')$ 的选取，
将在第 7.2 节证明。因此，在此假定命题 4.6.14，
并完成定理 4.6.1 的证明。先准备一些记号。

令 $J^\psi=\psi_*J$。相容殆复结构的空间可缩，
所以可取相容殆复结构的一族
$\{J_\rho\}_\rho=\{J_\rho\mid\rho\in[0,1]\}$，
满足
\[
 J_0=J^\psi(=\psi_*J),\qquad J_1=J'.
\]
微分同胚 $\psi\colon L\to L'$ 在 $L$ 与 $L'$
的光滑奇异链向量空间之间诱导同构。
令 $\mathcal X_1(L')$ 和 $C_1(L';\mathbb Q)$
分别为 $\mathcal X_1(L)$ 和 $C_1(L;\mathbb Q)$
在此同构下的像。用 $J^\psi$ 得到模空间
$\mathcal M_{k+1}^{\mathrm{main}}
((L',J^\psi),\beta;\vec P')$。
若 $P_i'=\psi(P_i)$，则作为带仓西结构的空间，
显然有
\begin{equation}
 \mathcal M_{k+1}^{\mathrm{main}}
 ((L',J^\psi),\beta;\vec P')
 \cong
 \mathcal M_{k+1}^{\mathrm{main}}
 ((L,J),\beta;\vec P).
\tag{4.6.4}\label{eq:4-6-4}
\end{equation}
因此，原模空间上的多值截面系 $\mathfrak s$
在左端模空间上诱导一套多值截面，记为
$\mathfrak s^\psi$。由此在
$C_1(L';\Lambda_{0,\mathrm{nov}})$
上得到滤过 $A_\infty$ 代数结构
$\mathfrak m_k^{J^\psi,\mathfrak s^\psi}$。
由 \eqref{eq:4-6-4}，$\psi$ 诱导滤过
$A_\infty$ 代数的典范同构（不只是同伦等价）
\[
 \psi_*\colon
 (C_1(L;\Lambda_{0,\mathrm{nov}}),
  \mathfrak m_k^{J,\mathfrak s})
 \longrightarrow
 (C_1(L';\Lambda_{0,\mathrm{nov}}),
  \mathfrak m_k^{J^\psi,\mathfrak s^\psi}).
\]

\begin{remark}{4.6.5}\label{rem:4-6-5}
由上面的讨论，在证明定理 4.6.1 时，
可以把 $L,J,\mathfrak s$ 换成
$L',J^\psi,\mathfrak s^\psi$，从而假定
$\psi=\mathrm{id}$。但我们不这样做，
因为下一小节将证明定理 4.6.1 所构造的
同伦等价不依赖于 $\psi$ 的辛同痕；
在那里需要考察 \eqref{eq:4-6-4}
所示同构的一个非常值族。
\end{remark}

为构造定理 4.6.1 中的 $\mathfrak f$，
我们使用另一个亏格为零的带边伪全纯曲线模空间。
在这里，各分支所用的相容殆复结构可以不同。
下面先定义这样的模空间。

设 $(\Sigma,\vec z)$ 是
$\mathcal M_{k+1}^{\mathrm{main}}$ 的一个元素，
$\{\Sigma_\alpha\mid\alpha\in\mathfrak A\}$
为 $\Sigma$ 的分支集。若 $\Sigma$ 有球面冒泡，
则把每棵球面冒泡树并入其所根植的圆盘分支；
特别地，每个 $\Sigma_\alpha$ 的边界都是 $S^1$。
我们在 $\mathfrak A$ 上如下定义偏序。
设 $\alpha_1,\alpha_2\in\mathfrak A$，
对应的分支为 $\Sigma_{\alpha_1},\Sigma_{\alpha_2}$；
令 $\Sigma_{\alpha_0}$ 为含第 $0$ 个标记点的分支。

\begin{definition}{4.6.6}\label{def:4-6-6}
若连接 $\Sigma_{\alpha_1}$ 与 $\Sigma_{\alpha_0}$
的每一条道路都同 $\Sigma_{\alpha_2}$ 相交，
则称 $\alpha_1\le\alpha_2$。
\end{definition}
容易看出，定义 4.6.6 确实给出
$\mathfrak A$ 上的偏序，其中 $\alpha_0$ 是最大元。
现在考察满足以下条件的系统
$((\Sigma,\vec z),(u_\alpha),(\rho_\alpha))$：
\begin{enumerate}[label=\textup{(4.6.7.\arabic*)},leftmargin=5.9em]
 \item
 $u_\alpha\colon(\Sigma_\alpha,\partial\Sigma_\alpha)
 \to(M',L')$ 是 $J_{\rho_\alpha}$-全纯的。
 \item
 $\rho_\alpha\in[0,1]$；若 $\alpha_1\le\alpha_2$，
 则 $\rho_{\alpha_1}\le\rho_{\alpha_2}$。
 \item
 $((\Sigma,\vec z),(u_\alpha))$
 在定义 2.1.24 的意义下是稳定的。
 \item
 $u_\alpha$ 的同调类为 $\beta(\alpha)$，
 且 $\sum_{\alpha\in\mathfrak A}\beta(\alpha)=\beta$。
 \item
 若 $z\in\Sigma_\alpha\cap\Sigma_{\alpha'}$，
 则 $u_\alpha(z)=u_{\alpha'}(z)$。
\end{enumerate}

\begin{figure}[htbp]
 \centering
 \setlength{\unitlength}{0.8mm}
 \begin{picture}(118,74)
  \put(66,27){\circle{27}}
  \put(53,48){\circle{22}}
  \put(45,16){\circle{19}}
  \put(39,3.5){\circle{8}}
  \put(36,48){\circle{11}}
  \put(55,63){\circle{8}}
  \put(61,26){$u_{\alpha_0}$}
  \put(48,47){$u_{\alpha_2}$}
  \put(41,15){$u_{\alpha_1}$}
  \put(36,2.5){$\scriptstyle u_{\alpha_3}$}
  \put(32,47){$\scriptstyle u_{\alpha_4}$}
  \put(52,62){$\scriptstyle u_{\alpha_5}$}
  \put(58.8,38.6){\makebox(0,0){$\times$}}
  \put(52.7,29.1){\makebox(0,0){$\times$}}
  \put(35.8,13.5){\makebox(0,0){$\times$}}
  \put(53.5,11.8){\makebox(0,0){$\times$}}
  \put(62.4,42.3){\makebox(0,0){$\times$}}
  \put(74,38){\makebox(0,0){$\times$}}
  \put(71,14.5){\makebox(0,0){$\times$}}
  \put(51.9,71.1){\circle*{0.55}}
  \put(54.9,71.5){\circle*{0.55}}
  \put(58.4,71.3){\circle*{0.55}}
  \put(61.0,70.3){\circle*{0.55}}
  \put(63.6,68.7){\circle*{0.55}}
  \put(65.8,66.5){\circle*{0.55}}
  \put(67.5,63.9){\circle*{0.55}}
  \put(68.5,61.0){\circle*{0.55}}
  \put(68.85,57.9){\circle*{0.55}}
  \put(68.5,54.8){\circle*{0.55}}
  \put(67.5,51.9){\circle*{0.55}}
  \qbezier(71.8,13.2)(72.125,12.55)(72.45,11.9)
  \qbezier(73.1,10.6)(73.425,9.95)(73.75,9.3)
  \qbezier(74.4,8.0)(74.725,7.35)(75.05,6.7)
  \qbezier(75.7,5.4)(76.025,4.75)(76.35,4.1)
  \put(79.5,27){\circle*{1}}
  \put(62.4,42.3){\circle*{1}}
  \put(74,38){\circle*{1}}
  \put(71,14.5){\circle*{1}}
  \put(80,26){$\longrightarrow z_0$}
  \put(64.5,41.8){$\scriptstyle z_2$}
  \put(75.5,39.5){$\scriptstyle z_1$}
  \put(73,11){$\scriptstyle z_k$}
  \put(3,53){$\rho_5\le\rho_2$}
  \put(3,42){$\rho_4\le\rho_2\le\rho_0$}
  \put(3,31){$\rho_3\le\rho_1\le\rho_0$}
  \put(80,2){$u_{\alpha_i}\text{ 是 }J_{\rho_i}\text{-全纯的}$}
 \end{picture}
 \par\smallskip
 \renewcommand{\thefigure}{4.6.1}
 \refstepcounter{figure}
 \textbf{\figurename~\thefigure}\label{fig:4-6-1}
\end{figure}

如果两个这样的系统之间存在自然意义下的同构，
就将它们视为等价。注意定义 4.6.6
类似于量子场论中的“时间有序积”。
\liftIndex{I173}

\begin{definition}{4.6.8}\label{def:4-6-8}
记
$\mathcal M_{k+1}^{\mathrm{main}}
(M',L',\{J_\rho\}_\rho:\beta;\operatorname{top}(\rho))$
为满足 (4.6.7) 的系统
$((\Sigma,\vec z),(u_\alpha),(\rho_\alpha))$
的所有等价类所成的集合。
\end{definition}

此处及下文，记号
$\mathcal M_{k+1}^{\mathrm{main}}
(M',L',\{J_\rho\}_\rho:\beta;\operatorname{top}(\rho))$
中的 $\operatorname{top}(\rho)$
表示对参数 $\rho$ 取上文意义的时间有序积。
以后考察由若干变量参数化的殆复结构族所定义的
模空间时，将按需要取这样的时间有序积，
或者在每一个时间上取纤维积（逐时积）。
\liftIndex{I174}
后一种情形用 $\operatorname{twp}(s)$ 标明。
这两种模空间分别对应于证明弗洛尔上同调
不依赖于各种选择的两种方法。
参见式 (2.3.13) 后的讨论及第 7.2.14 小节。

与往常一样，模空间
$\mathcal M_{k+1}^{\mathrm{main}}
(M',L',\{J_\rho\}_\rho:\beta;\operatorname{top}(\rho))$
有带边界的仓西结构。
其中一列元素
$((\Sigma^{(i)},\vec z^{\,(i)}),
(u_\alpha^{(i)}),(\rho_\alpha^{(i)}))$
若发生下列情况之一，便趋于仓西结构意义下的边界：
\begin{enumerate}[label=\textup{(4.6.9.\arabic*)},leftmargin=5.9em]
 \item
 某个分支 $\Sigma_\alpha^{(i)}$ 裂成两个分支。
 \item
 对相交的分支
 $\Sigma_\alpha^{(i)}\cap\Sigma_{\alpha'}^{(i)}
 \ne\varnothing$，
 有
 $\lim_{i\to\infty}\rho_\alpha^{(i)}
 =\lim_{i\to\infty}\rho_{\alpha'}^{(i)}$。
 \item
 $\lim_{i\to\infty}\rho_\alpha^{(i)}=0$。
 \item
 $\lim_{i\to\infty}\rho_\alpha^{(i)}=1$。
\end{enumerate}
见图 4.6.2--4.6.5。
与 (4.6.9.1)--(4.6.9.4) 对应的边界分支，
每个都可写成模空间的纤维积。
带仓西结构空间的纤维积定义见第 A1.2 节。

\begin{figure}[htbp]
 \centering
 \setlength{\unitlength}{0.8mm}
 \begin{picture}(133,70)
  \put(9,61){\circle{8}}
  \qbezier(20.8,47.2)(22.6,55)(14.8,56.8)
  \qbezier(14.8,56.8)(7,58.6)(5.2,50.8)
  \qbezier(5.2,50.8)(3.4,43)(11.2,41.2)
  \qbezier(11.2,41.2)(12.6,40.9)(14,41.1)
  \qbezier(20.8,47.2)(17.9,38.9)(26.5,41)
  \qbezier(14,41.1)(22.6,43.1)(19.7,34.8)
  \qbezier(19.7,34.8)(17.9,27)(25.7,25.2)
  \qbezier(25.7,25.2)(33.5,23.4)(35.3,31.2)
  \qbezier(35.3,31.2)(38,37.2)(29.3,40.8)
  \qbezier(29.3,40.8)(27.9,41.1)(26.5,41)
  \put(37,20){\circle{16}}
  \put(36,6){\circle{10}}
  \put(45.5,20){\vector(1,0){2.5}}
  \put(49,19){$ev_0$}
  \put(42,38.5){$\Longrightarrow$}
  \put(45.5,29){$i\to\infty$}
  \put(86.5,54){\circle{8}}
  \put(89,42){\circle{16}}
  \put(100.5,30){\circle{16}}
  \put(112.5,19){\circle{16}}
  \put(109,6){\circle{10}}
  \put(121,19){\vector(1,0){2.5}}
  \put(124.5,18.5){$ev_0$}
 \end{picture}
 \par\smallskip
 \renewcommand{\thefigure}{4.6.2}
 \refstepcounter{figure}
 \textbf{\figurename~\thefigure}\label{fig:4-6-2}
 \textup{(4.6.9.1)}
\end{figure}

\begin{figure}[htbp]
 \centering
 \setlength{\unitlength}{0.8mm}
 \begin{picture}(100,58)
  \put(28,48){\circle{10}}
  \put(31,35){\circle{16}}
  \put(45,25){\circle{16}}
  \put(62,16){\circle{19}}
  \put(70,16){$\longrightarrow ev_0$}
  \put(25,34){$\scriptstyle\rho_\alpha^{(i)}$}
  \put(39,24){$\scriptstyle\rho_{\alpha'}^{(i)}$}
  \put(22,0){$\rho_\alpha^{(i)}
   -\rho_{\alpha'}^{(i)}\longrightarrow0$}
 \end{picture}
 \par\smallskip
 \renewcommand{\thefigure}{4.6.3}
 \refstepcounter{figure}
 \textbf{\figurename~\thefigure}\label{fig:4-6-3}
 \textup{(4.6.9.2)}
\end{figure}

\begin{figure}[htbp]
 \centering
 \setlength{\unitlength}{0.8mm}
 \begin{picture}(105,65)
  \put(76,19){\circle{26}}
  \put(57,31){\circle{21}}
  \put(36,38){\circle{19}}
  \put(34,53){\circle{12}}
  \put(25,28){\circle{12}}
  \put(65,5){\circle{13}}
  \put(84,29){\circle*{1}}
  \put(64,39){\circle*{1}}
  \put(38,58){\circle*{1}}
  \put(28,53){\circle*{1}}
  \put(19,28){\circle*{1}}
  \put(36,28.5){\circle*{1}}
  \put(61,0){\circle*{1}}
  \put(76,6){\circle*{1}}
  \put(72,18){$\rho_0$}
  \put(53,30){$\rho_2$}
  \put(32,37){$\rho_3$}
  \put(42,52){$\rho_5=0$}
  \put(21,27){$\rho_4$}
  \put(61,4){$\rho_1$}
  \put(83,36){$P_1$}
  \put(60,43){$P_2$}
  \put(40,61){$P_3$}
  \put(15,52){$P_4$}
  \put(13,30){$P_5$}
  \put(31,16){$P_6$}
  \put(51,0){$P_7$}
  \put(73,0){$P_8$}
  \put(88,20){$\longrightarrow ev_0$}
 \end{picture}
 \par\smallskip
 \renewcommand{\thefigure}{4.6.4}
 \refstepcounter{figure}
 \textbf{\figurename~\thefigure}\label{fig:4-6-4}
 \textup{(4.6.9.3)}
\end{figure}

\begin{figure}[htbp]
 \centering
 \setlength{\unitlength}{0.8mm}
 \begin{picture}(100,57)
  \put(68,22){\circle{28}}
  \put(50,35){\circle{21}}
  \put(42,51){\circle{11}}
  \put(58,5){\circle{13}}
  \put(76,33){\circle*{1}}
  \put(57,43){\circle*{1}}
  \put(40,35){\circle*{1}}
  \put(54,22){\circle*{1}}
  \put(54,0){\circle*{1}}
  \put(68,8){\circle*{1}}
  \put(63,21){$\rho_0=1$}
  \put(46,34){$\rho_2$}
  \put(39,50){$\rho_3$}
  \put(54,4){$\rho_1$}
  \put(78,37){$P_1$}
  \put(58,46){$P_2$}
  \put(32,36){$P_3$}
  \put(51,20){$P_4$}
  \put(45,0){$P_5$}
  \put(70,3){$P_6$}
  \put(82,22){$\longrightarrow ev_0$}
 \end{picture}
 \par\smallskip
 \renewcommand{\thefigure}{4.6.5}
 \refstepcounter{figure}
 \textbf{\figurename~\thefigure}\label{fig:4-6-5}
 \textup{(4.6.9.4)}
\end{figure}

注意 (4.6.9.1) 与 (4.6.9.2) 相抵消，
故只需考察 (4.6.9.3) 与 (4.6.9.4)。
描述 (4.6.9.3) 的模空间，是若干
$\mathcal M_{\ell+1}^{\mathrm{main}}
(M',L',\{J_\rho\}_\rho:\beta';
\operatorname{top}(\rho))$（$\beta'\le\beta$）
和
$\mathcal M_{\ell'+1}^{\mathrm{main}}
((L',J^\psi),\beta'')$
的纤维积，因为 $J_0=J^\psi$。
类似地，描述 (4.6.9.4) 的模空间，
是前一种模空间（$\beta'\le\beta$）
和
$\mathcal M_{\ell'+1}^{\mathrm{main}}
((L',J'),\beta'')$
的纤维积。我们将在证明末尾回到这一点；
它是定理 4.6.1 证明的主要思路。

定义评价映射
\[
 ev=(ev_1,\ldots,ev_k,ev_0)\colon
 \mathcal M_{k+1}^{\mathrm{main}}
 (M',L',\{J_\rho\}_\rho:\beta;\operatorname{top}(\rho))
 \longrightarrow L^k\times L'
\]
如下：
\[
 \begin{cases}
 ev_0((\Sigma,\vec z),(u_\alpha),(\rho_\alpha))
 =u(z_0),\\
 ev_i((\Sigma,\vec z),(u_\alpha),(\rho_\alpha))
 =\psi^{-1}(u(z_i)),
 \quad i=1,\ldots,k.
 \end{cases}
\]

\begin{definition}{4.6.10}\label{def:4-6-10}
对 $P_i\in\mathcal X_1(L)$，定义
\[
 \begin{aligned}
 &\mathcal M_{k+1}^{\mathrm{main}}
 (M',L',\{J_\rho\}_\rho:\beta;\operatorname{top}(\rho);\vec P)\\
 &\quad=
 \mathcal M_{k+1}^{\mathrm{main}}
 (M',L',\{J_\rho\}_\rho:\beta;\operatorname{top}(\rho))
 {}_{(ev_1,\ldots,ev_k)}\!\times
 \left(\prod_{i=1}^kP_i\right).
 \end{aligned}
\]
右端是以 $L^k$ 为底的纤维积。
\end{definition}

令 $\mathfrak s^{\operatorname{top}(\rho)}$
为
$\mathcal M_{k+1}^{\mathrm{main}}
(M',L',\{J_\rho\}_\rho:\beta;\operatorname{top}(\rho);\vec P)$
上的一套多值截面。
它应满足的性质将在下面的命题 4.6.14 说明。
置
\begin{equation}
 \begin{aligned}
 &\mathcal M_{k+1}^{\mathrm{main}}
 (M',L',\{J_\rho\}_\rho:\beta;\operatorname{top}(\rho);\vec P)
 ^{\mathfrak s^{\operatorname{top}(\rho)}}\\
 &\qquad=
 ((\mathfrak s^{\operatorname{top}(\rho)})^{-1}(0),ev_0).
 \end{aligned}
\tag{4.6.11}\label{eq:4-6-11}
\end{equation}
这里通过适当三角剖分，把右端看成奇异链；
参见定义 A1.28。关于此空间的定向，参见第 8.9 节。
现在给出定义。

\begin{definition}{4.6.12}\label{def:4-6-12}
\begin{enumerate}[label=\textup{(4.6.13.\arabic*)},leftmargin=6.1em]
 \item
 $\mathfrak f_{1,\beta_0}(P)=\psi(P)$。
 \item
 当 $(k,\beta)\ne(1,\beta_0)$ 时，
 \[
 \mathfrak f_{k,\beta}(P_1,\ldots,P_k)
 =
 \mathcal M_{k+1}^{\mathrm{main}}
 (M',L',\{J_\rho\}_\rho:\beta;\operatorname{top}(\rho);\vec P)
 ^{\mathfrak s^{\operatorname{top}(\rho)}}.
 \]
 \item
 \[
 \mathfrak f_k(P_1,\ldots,P_k)
 =
 \sum_\beta
 \mathfrak f_{k,\beta}(P_1,\ldots,P_k)
 \otimes T^{E(\beta)}e^{\mu_L(\beta)/2}.
 \]
\end{enumerate}
\end{definition}

由格罗莫夫紧性，对任意 $E_0$，
使
$\mathcal M_{k+1}^{\mathrm{main}}
(M',L',\{J_\rho\}_\rho:\beta;\operatorname{top}(\rho))\ne\varnothing$
且 $E(\beta)<E_0$ 的 $\beta$ 只有有限多个。
因此，$\mathfrak f_k(P_1,\ldots,P_k)$ 良定义，
而 $\mathfrak f$ 是有隙的。

为证明定理 4.6.1，还须说明 (4.6.13) 的右端
究竟落在哪个空间。这涉及可数集
$\mathcal X_3(L')$ 的选择；
我们将在第 7.2.9 小节详细证明这一点
（见命题 7.2.100）。现在可陈述
$\mathfrak s^{\operatorname{top}(\rho)}$
应满足的性质。

\begin{proposition}{4.6.14}\label{prop:4-6-14}
可以找到 $\mathcal X_3(L')$、
$\mathfrak s''$ 和
$\mathfrak s^{\operatorname{top}(\rho)}$，
使以下性质成立：
\begin{enumerate}[label=\textup{(4.6.15.\arabic*)},leftmargin=6.1em]
 \item
 (4.6.2.1)、(4.6.2.2)、(4.6.2.3) 成立。
 \item
 对任意 $P_i\in\mathcal X_1(L)$，
 由 (4.6.13) 定义的有理数系数链
 $\mathfrak f_{k,\beta}(P_1,\ldots,P_k)$
 属于 $C_3(L';\mathbb Q)$。
 换言之，
 $\mathfrak f_k(P_1,\ldots,P_k)
 \in C_3(L';\Lambda_{0,\mathrm{nov}})$。
 尤其对每个 $P\in\mathcal X_1(L)$，
 $\psi(P)\in C_3(L';\Lambda_{0,\mathrm{nov}})$。
 \item
 在与 (4.6.9.3) 对应的
 $\mathcal M_{k+1}^{\mathrm{main}}
 (M',L',\{J_\rho\}_\rho:\beta;\operatorname{top}(\rho);\vec P)$
 的边界上，
 $\mathfrak s^{\operatorname{top}(\rho)}$
 与 $\mathfrak s^\psi$ 相容。
 \item
 在与 (4.6.9.4) 对应的同一模空间边界上，
 $\mathfrak s^{\operatorname{top}(\rho)}$
 与 $\mathfrak s''$ 相容。
\end{enumerate}
\end{proposition}

命题 4.6.14 将在第 7.2.9 小节证明。
准确地说，在那里先固定 $E_0$，
只对满足 $\omega[\beta]\le E_0$
的 $\beta$ 证明命题 4.6.14。
这对任意 $n,K$ 定义一个滤过 $A_{n,K}$ 同态
（意义见第 7.2.6 小节）。
随后将在第 7.2.11 小节运用同调代数
构造滤过 $A_\infty$ 同态；
另参见注 3.5.13。

这里假定命题 4.6.14，完成定理 4.6.1
的证明。为此只须验证 $\mathfrak f$
是滤过 $A_\infty$ 同态。
令 $P_j\in\mathcal X_1(L)$，写作
$P_j=(|P_j|,f_j)$，其中 $|P_j|$ 是单形，
$f_j\colon|P_j|\to L$ 是光滑映射。
考察链
$\mathfrak f_{k,\beta}(P_1,\ldots,P_k)$ 的边界。
设
$((\Sigma^{(i)},\vec z^{\,(i)}),
(u_\alpha^{(i)}),(\rho_\alpha^{(i)}))$
是一列元素，$x_j^{(i)}\in|P_j|$，
且对 $j=1,\ldots,k$ 有
\[
 ev_j((\Sigma^{(i)},\vec z^{\,(i)}),
 (u_\alpha^{(i)}),(\rho_\alpha^{(i)}))
 =f_j(x_j^{(i)}).
\]
于是得到模空间
$\mathcal M_{k+1}^{\mathrm{main}}
(M',L',\{J_\rho\}_\rho:\beta;\operatorname{top}(\rho);\vec P)$
中的一列元素
\[
 ((\Sigma^{(i)},\vec z^{\,(i)}),
 (u_\alpha^{(i)}),(\rho_\alpha^{(i)});
 x_1^{(i)},\ldots,x_k^{(i)}).
\]
当下列情形之一发生时，这列元素趋于边界：
\begin{enumerate}[label=\textup{(4.6.16.\arabic*)},leftmargin=6.1em]
 \item 发生 (4.6.9.1)。
 \item 发生 (4.6.9.2)。
 \item 发生 (4.6.9.3)。
 \item 发生 (4.6.9.4)。
 \item 某个 $x_j^{(i)}$ 趋于 $P_j$ 的边界。
\end{enumerate}

(4.6.16.1) 与 (4.6.16.2) 如前所述相抵消。
(4.6.16.3) 对应于如下求和；求和须包含几何上允许的
零类项。\translatornote{原书 (4.6.17)、(4.6.18) 的求和
下限分别写有 $\beta_1,\beta_2\ne0$ 与
$\beta_1,\beta_{2,i}\ne0$。若逐字排除零类，
则图 4.6.5 对应式中明确出现的
$\mathfrak f_{1,\beta_0}(P)=\psi(P)$
无法计入 (4.6.18)；(4.6.17) 也会排除
零类的经典乘法项。两式在此按完整
$A_\infty$ 同态关系求和，保留稳定性等原有适用条件。}
\begin{equation}
 \begin{aligned}
 \sum_{\substack{1\le\ell\le m\le k\\
                  \beta=\beta_1+\beta_2}}
 &\mathfrak f_{k+\ell-m,\beta_1}
 \bigl(P_1,\ldots,P_{\ell-1},\\
 &\quad\mathfrak m^{J,\mathfrak s}_{m-\ell+1,\beta_2}
 (P_\ell,\ldots,P_m),
 P_{m+1},\ldots,P_k\bigr).
 \end{aligned}
\tag{4.6.17}\label{eq:4-6-17}
\end{equation}
例如，图 4.6.4 对应于
\[
 \mathfrak f_7\bigl(
 P_1,P_2,\mathfrak m_2(P_3,P_4),
 P_5,P_6,P_7,P_8\bigr).
\]

(4.6.16.4) 对应于求和
\begin{equation}
 \begin{aligned}
 \sum_{\substack{m,\ 1=\ell_1\le\ell_2\le
                   \cdots\le\ell_m=k+1\\
                  \beta=\beta_1+
                   \sum_{i=1}^{m-1}\beta_{2,i}}}
 &\mathfrak m^{J',\mathfrak s'}_{m-1,\beta_1}
 \Bigl(
 \mathfrak f_{\ell_2-\ell_1,\beta_{2,1}}
 (P_1,\ldots,P_{\ell_2-1}),\\
 &\qquad\ldots,\,
 \mathfrak f_{\ell_{i+1}-\ell_i,\beta_{2,i}}
 (P_{\ell_i},\ldots,P_{\ell_{i+1}-1}),
 \ldots,\\
 &\qquad
 \mathfrak f_{\ell_m-\ell_{m-1},\beta_{2,m-1}}
 (P_{\ell_{m-1}},\ldots,P_{\ell_m-1})
 \Bigr).
 \end{aligned}
\tag{4.6.18}\label{eq:4-6-18}
\end{equation}
例如，图 4.6.5 对应于
\[
 \mathfrak m_5\bigl(
 \psi(P_1),\mathfrak f_2(P_2,P_3),
 \psi(P_4),\mathfrak f_1(P_5),\psi(P_6)\bigr).
\]
这里 $\psi(P)=\mathfrak f_{1,\beta_0}(P)$。

(4.6.16.5) 对应于求和
\begin{equation}
 \sum_j\mathfrak f_{k,\beta}
 (P_1,\ldots,\overline{\mathfrak m}_1(P_j),
  \ldots,P_k).
\tag{4.6.19}\label{eq:4-6-19}
\end{equation}
因此，
$\overline{\mathfrak m}_1
(\mathfrak f_{k,\beta}(P_1,\ldots,P_k))$
等于 \eqref{eq:4-6-17}、
\eqref{eq:4-6-18} 与 \eqref{eq:4-6-19}
之和。这里暂不讨论符号；符号问题见第 8.9 节。
这就说明 $\mathfrak f$ 是滤过 $A_\infty$ 同态。

\begin{remark}{4.6.20}\label{rem:4-6-20}
在上述论证中，需要用多值截面扰动
定义 4.6.10 的模空间，并使用
\eqref{eq:4-6-11} 的扰动后模空间。
不过，命题 4.6.14 所述多值截面的相容性，
正是为了使上述论证在计入
$\mathfrak s^{\operatorname{top}(\rho)}$
等所诱导的扰动后仍然成立。
\end{remark}
